\documentclass[addpoints, 11pt]{exam}
\usepackage[a4paper, left=20mm, right=20mm, top=26mm, bottom=30mm, headsep=6mm, footskip=12mm, headheight=15pt]{geometry}
\usepackage[T1]{fontenc}
\usepackage[utf8]{inputenc}
\usepackage{lmodern}
\usepackage{textcomp}
\usepackage[french]{babel}
\usepackage{xcolor}
\usepackage{listings}
\usepackage{graphicx}
\usepackage{array}
\usepackage{amsmath}
\usepackage{eurosym}
\usepackage{needspace}
\usepackage[most]{tcolorbox}
\usepackage{tikz}
\usetikzlibrary{positioning,arrows.meta}


% --- Code Python : même style que le sujet, en police à chasse fixe ---
\definecolor{codegreen}{rgb}{0,0.6,0}
\definecolor{codegray}{rgb}{0.5,0.5,0.5}
\definecolor{codepurple}{rgb}{0.58,0,0.82}
\definecolor{backcolour}{rgb}{0.95,0.95,0.92}
\lstdefinestyle{mystyle}{
    backgroundcolor=\color{backcolour},
    commentstyle=\color{codegreen},
    keywordstyle=\color{magenta},
    numberstyle=\tiny\color{codegray},
    stringstyle=\color{codepurple},
    basicstyle=\footnotesize\ttfamily,
    columns=fullflexible,
    breakatwhitespace=false,
    breaklines=true,
    postbreak=\mbox{\textcolor{codegray}{$\hookrightarrow$}\space},
    captionpos=b,
    keepspaces=true,
    numbers=left,
    numbersep=5pt,
    xleftmargin=1.5em,
    showspaces=false,
    showstringspaces=false,
    showtabs=false,
    tabsize=4,
    upquote=true
}
\lstset{style=mystyle, language=Python, extendedchars=true,
  literate={à}{{\`a}}1 {â}{{\^a}}1 {ä}{{\"a}}1 {ç}{{\c{c}}}1
           {é}{{\'e}}1 {è}{{\`e}}1 {ê}{{\^e}}1 {ë}{{\"e}}1
           {î}{{\^i}}1 {ï}{{\"i}}1 {ô}{{\^o}}1 {ö}{{\"o}}1
           {ù}{{\`u}}1 {û}{{\^u}}1 {ü}{{\"u}}1 {œ}{{\oe}}2
           {É}{{\'E}}1 {È}{{\`E}}1 {À}{{\`A}}1 {Ç}{{\c{C}}}1
           {°}{{\textdegree}}1 {€}{{\euro}}1 {…}{{\dots}}1
           {«}{{\guillemotleft}}1 {»}{{\guillemotright}}1 {·}{{\textperiodcentered}}1}
% Terminal (commandes Linux) : sans couleur de mots-clés
\lstdefinestyle{terminal}{style=mystyle, language={}, numbers=none, xleftmargin=0.5em}
% Sortie d'un programme (ce qui s'affiche à l'écran)
\lstdefinestyle{sortie}{style=mystyle, language={}, numbers=none, backgroundcolor=\color{white}, frame=single, rulecolor=\color{black!60}, xleftmargin=0.5em}

% --- Boîtes du corrigé (lisibles en noir et blanc) ---
\newtcolorbox{reponse}[1][Réponse attendue]{enhanced, breakable, colback=white, colframe=black, boxrule=1.2pt, arc=0pt,
  fonttitle=\bfseries, coltitle=white, colbacktitle=black, title=#1, left=2mm, right=2mm, top=1mm, bottom=1mm, before skip=6pt, after skip=6pt}
\newtcolorbox{erreur}[1][Erreur fréquente]{enhanced, breakable, colback=black!6, colframe=black!55, boxrule=0.6pt, arc=2pt,
  borderline west={3pt}{0pt}{black!70}, fonttitle=\bfseries, coltitle=black, colbacktitle=black!15, title=#1, left=2mm, right=2mm, top=1mm, bottom=1mm, before skip=6pt, after skip=6pt}
\newtcolorbox{astuce}[1][À retenir]{enhanced, breakable, colback=white, colframe=black!55, boxrule=0.6pt, arc=2pt,
  fonttitle=\bfseries, coltitle=black, colbacktitle=black!8, title=#1, left=2mm, right=2mm, top=1mm, bottom=1mm, before skip=6pt, after skip=6pt}
\newtcolorbox{remarque}[1][Remarque sur l'énoncé]{enhanced, breakable, colback=white, colframe=black!70, boxrule=0.6pt, arc=0pt,
  borderline={0.4pt}{1.2pt}{black!70, dashed}, fonttitle=\bfseries, coltitle=black, colbacktitle=white, title=#1, left=2mm, right=2mm, top=1mm, bottom=1mm, before skip=6pt, after skip=6pt}
\newcommand{\enonce}[1]{\textit{#1}\par\smallskip}
\newcommand{\methode}{\par\smallskip\noindent\textbf{Méthode.}\ }
\newcommand{\explication}{\par\smallskip\noindent\textbf{Explication.}\ }
\newcommand{\code}[1]{\texttt{#1}}
% Coupure de ligne possible après un tiret bas (noms de variables longs)
\DeclareRobustCommand{\_}{\textunderscore\allowbreak}
% Un code court ne se coupe pas en bas de page
\AddToHook{env/lstlisting/before}{\needspace{10\baselineskip}}
\setlength{\parindent}{0pt}
\setlength{\parskip}{3pt}
\renewcommand{\arraystretch}{1.25}

\title{Corrigé --- CC 1re NSI - Fonctions (2024)}
\author{Babash - NSI Tani Malandi et Johanne Bakalara}
\date{Novembre 2024, durée : 1 h 30}

\begin{document}
\pagestyle{headandfoot}
\runningheader{}{}{\footnotesize Corrigé --- Fonctions (2024)}
\firstpageheader{}{}{}
\runningfooter{}{\footnotesize Babash - NSI Tani Malandi et Johanne Bakalara \textperiodcentered{} CC BY-SA 4.0 \textperiodcentered{} page \thepage}{}
\firstpagefooter{}{\footnotesize Babash - NSI Tani Malandi et Johanne Bakalara \textperiodcentered{} CC BY-SA 4.0 \textperiodcentered{} page \thepage}{}

\maketitle
\thispagestyle{headandfoot}

\begin{astuce}[Comment lire ce corrigé]
Pour chaque question : la question est rappelée en italique, puis la \textbf{réponse attendue} est encadrée en noir, puis l'\textbf{explication} montre la méthode, étape par étape.
Les encadrés gris « Erreur fréquente » signalent les pièges les plus courants. Essaie d'abord de refaire la question seul·e, puis compare ta démarche avec celle du corrigé.
\end{astuce}

\section*{Barème}
\begin{center}
  \gradetable[h]
\end{center}

\begin{questions}

% =====================================================================
\section{Question de cours}

\question[1] \enonce{Une fonction est-elle une structure de contrôle : vrai ou faux ?}
\begin{reponse}
\textbf{Vrai.} Dans la convention du cours, une structure de contrôle est tout ce qui prend un \code{:} en Python : \code{def} en fait partie, comme \code{if}, \code{elif}, \code{else}, \code{for} et \code{while}.
\end{reponse}
\explication\ Une \textbf{structure de contrôle} est une instruction qui ouvre un \textbf{bloc} : elle se termine par deux-points (\code{:}) et son contenu est indenté. La liste des mots concernés : \code{if}, \code{elif}, \code{else}, \code{for}, \code{while}, \code{def}, et plus largement \code{class}, \code{try}, \code{except}, \code{finally}, \code{with}, \code{match} et \code{case} quand on les utilise. Une \textbf{fonction} (\code{def}) est donc bien une structure de contrôle : c'est un bloc d'instructions auquel on donne un nom, pour le réutiliser en l'appelant. Elle n'est ni une instruction conditionnelle ni une boucle, mais comme elles elle ouvre un bloc après un \code{:}.
\begin{remarque}
Certains manuels réservent l'expression aux instructions conditionnelles et aux boucles. La réponse attendue ici suit la convention du cours : « tout ce qui prend un \code{:} ».
\end{remarque}

\question[3] \enonce{Soit la fonction \code{somme\_entre\_deux\_nombres}.}
%%% Bloc commun aux corrigés de fonctions (2024 et 2025) : questions sur somme_entre_deux_nombres
\begin{lstlisting}
def somme_entre_deux_nombres(a, b):
    somme = 0
    for i in range(a, b):
        somme += i
    return somme
\end{lstlisting}
\begin{parts}
\part \enonce{Quels sont les paramètres de la fonction ?}
\begin{reponse}
\code{a} et \code{b}.
\end{reponse}
Les \textbf{paramètres} sont les noms écrits entre parenthèses dans l'en-tête \code{def ...(a, b):}. Ce sont des variables qui reçoivent les valeurs données lors de l'appel.

\part \enonce{Quelles sont les variables locales de la fonction ?}
\begin{reponse}
\code{somme} et \code{i} (ainsi que les paramètres \code{a} et \code{b}, qui sont eux aussi locaux).
\end{reponse}
Une variable \textbf{locale} est créée à l'intérieur de la fonction et n'existe que pendant son exécution : après l'appel, \code{somme} et \code{i} ne sont plus accessibles depuis le programme principal. \code{somme} est créée ligne 2, \code{i} est créée par la boucle \code{for} ligne 3.

\part \enonce{Quelles sont les valeurs de retour de la fonction ?}
\begin{reponse}
Une seule valeur de retour : \code{somme}, c'est-à-dire la somme des entiers de \code{a} (inclus) à \code{b} (exclu).
\end{reponse}
La valeur de retour est ce qui suit le mot-clé \code{return}. Attention : \code{range(a, b)} s'arrête \textbf{avant} \code{b} ; \code{b} n'est pas ajouté.
\end{parts}


\question[2] \enonce{Écrire un appel de \code{somme\_entre\_deux\_nombres} pour \code{a=10} et \code{b=15}, en stockant les valeurs de retour dans des variables.}
\begin{reponse}
\begin{lstlisting}
resultat = somme_entre_deux_nombres(10, 15)
\end{lstlisting}
\code{resultat} vaut \textbf{60}.
\end{reponse}
\explication
\begin{enumerate}
  \item On écrit le nom de la fonction, puis les \textbf{arguments} entre parenthèses, dans l'ordre des paramètres : 10 va dans \code{a}, 15 va dans \code{b}.
  \item La fonction renvoie \textbf{une} valeur : on la range dans \textbf{une} variable avec \code{=}. Sans variable, la valeur renvoyée serait perdue.
\end{enumerate}
\textbf{Trace de l'appel :}
\begin{center}
\begin{tabular}{|c|c|c|}
\hline
Tour & \code{i} & \code{somme} \\ \hline
avant & & 0 \\
1 & 10 & 10 \\
2 & 11 & 21 \\
3 & 12 & 33 \\
4 & 13 & 46 \\
5 & 14 & 60 \\ \hline
\end{tabular}
\end{center}
\code{range(10, 15)} donne 10, 11, 12, 13, 14 : 15 est exclu.
\begin{erreur}
Écrire \code{def somme\_entre\_deux\_nombres(10, 15)} : le mot \code{def} sert à \textbf{définir}, pas à appeler. Autre erreur : appeler sans ranger le résultat (\code{somme\_entre\_deux\_nombres(10, 15)} seul) : la valeur 60 est calculée puis perdue.
\end{erreur}

\question[4] \enonce{Associer à chaque numéro du schéma un terme : appeler, arguments, valeur de retour, corps, en-tête, définir, paramètre, nom.}
\begingroup
\tikzset{every picture/.append style={scale=0.78, transform shape}}
\begin{center}
\begin{tikzpicture}[font=\small,
  num/.style={draw,minimum width=0.7cm,minimum height=0.55cm,inner sep=1pt},
  lab/.style={anchor=west,inner sep=2pt},
  box/.style={draw=cyan!60!blue,rounded corners=6pt,very thick,inner sep=7pt}]
\node[box,anchor=center,align=left,font=\ttfamily\small] (g) at (0,0) {%
\textcolor{blue!70!black}{def} nbrCaractere (ch):\\
\hspace*{1.4em}i=0\\
\hspace*{1.4em}\textcolor{orange!80!black}{for} j \textcolor{orange!80!black}{in} ch:\\
\hspace*{2.8em}i+=1\\
\hspace*{1.4em}\textcolor{brown!80!black}{return} (i)};
\node[box,font=\ttfamily\small] (d) at (8,0) {nbrCaractere ('Alex')};
\node[num] (n1) at (-3,3.3) {1}; \node[lab] at (n1.east) {une fonction};
\node[num] (n2) at (5.6,1.6) {2}; \node[lab] at (n2.east) {une fonction};
\node[num] (n6) at (-1.6,2.3) {6}; \node[lab] at (n6.east) {de la fonction};
\node[num] (n8) at (2.6,2.3) {8};
\node[num] (n3) at (-6.6,1.0) {3}; \node[lab] (t3) at (n3.east) {de la fonction};
\node[num] (n4) at (-6.6,0) {4}; \node[lab] (t4) at (n4.east) {de la fonction};
\node[num] (n5) at (-6.6,-1.0) {5};
\node[num] (n7) at (8.3,-1.8) {7}; \node[lab] at (n7.east) {de la fonction};
\draw[->,thick] (t3.east) -- (-2.5,1.0);
\draw[->,thick] (t4.east) -- (-2.5,0);
\draw[->,thick] (n5.east) -- (-2.5,-1.0);
\draw[->,thick] (n6.south) -- (-0.6,1.15);
\draw[->,thick] (n8.south) -- (1.5,1.15);
\draw[->,thick] (n7.north) -- (8.3,-0.45);
\end{tikzpicture}
\end{center}

\endgroup
\begin{reponse}
\begin{tabular}{cl@{\qquad}cl}
1 & \textbf{définir} une fonction & 5 & \textbf{valeur de retour} \\
2 & \textbf{appeler} une fonction & 6 & \textbf{nom} de la fonction \\
3 & \textbf{en-tête} de la fonction & 7 & \textbf{arguments} de la fonction \\
4 & \textbf{corps} de la fonction & 8 & \textbf{paramètre} \\
\end{tabular}
\end{reponse}
\methode\ Commence par les mots que tu reconnais sûrement, puis termine par élimination.
\begin{enumerate}
  \item Le schéma a deux parties. À gauche, le bloc \code{def} est la \textbf{définition}.\\ À droite, la ligne \code{nbrCaractere('Alex')} est l'\textbf{appel}. Le 1 est au-dessus de la définition : « \textbf{définir} une fonction » ; le 2 est au-dessus de l'appel : « \textbf{appeler} une fonction ».
  \item Le 6 pointe vers le mot \code{nbrCaractere} : c'est le \textbf{nom}. Le 8 pointe vers \code{(ch)} : c'est le \textbf{paramètre}.
  \item Le 3 pointe vers la première ligne, celle du \code{def} : l'\textbf{en-tête}. Le 4 pointe vers les lignes décalées en dessous : le \textbf{corps}.
  \item Le 5 pointe vers la ligne \code{return (i)} : ce qui est renvoyé est la \textbf{valeur de retour}.
  \item Le 7 pointe vers \code{'Alex'}, la valeur donnée lors de l'appel : c'est l'\textbf{argument}.
\end{enumerate}
\begin{erreur}
Confondre \textbf{paramètre} et \textbf{argument}. Le paramètre (\code{ch}) est le nom écrit dans la définition ; l'argument (\code{'Alex'}) est la valeur concrète donnée à l'appel. Lors de l'appel, \code{ch} reçoit \code{'Alex'}, et la fonction renvoie 4 (le nombre de caractères).
\end{erreur}


% =====================================================================
\section{Compréhension de code}

\question[3] \enonce{Dans le code du jeu « devine le nombre », entourer les définitions des fonctions (en entier) et souligner les appels.}
\begin{reponse}
\textbf{Définitions à entourer} (chacune en entier, de la ligne \code{def} à la dernière ligne indentée) :
\begin{itemize}
  \item \code{demander\_nombre} : lignes 3 à 5 ;
  \item \code{aide} : lignes 7 à 11.
\end{itemize}
\textbf{Appels à souligner} :
\begin{itemize}
  \item ligne 20 : \code{demander\_nombre()} ;
  \item ligne 21 : \code{aide(rep, nbHasard)}.
\end{itemize}
On accepte aussi de souligner les appels de fonctions prédéfinies : \code{int(...)} et \code{input(...)} (ligne 4), \code{print(...)} (lignes 9, 11, 24, 26), \code{random.randint(1, 100)} (ligne 15).
\end{reponse}
\textbf{Le code annoté :}
\begin{lstlisting}[escapechar=@, numbers=none]
import random

@\fbox{\begin{minipage}{0.75\linewidth}\ttfamily\footnotesize\textcolor{magenta}{def} demander\_nombre():\\
\hspace*{2em}rep = int(input("quel est votre choix ?: "))\\
\hspace*{2em}\textcolor{magenta}{return} rep\end{minipage}}@

@\fbox{\begin{minipage}{0.75\linewidth}\ttfamily\footnotesize\textcolor{magenta}{def} aide(rep, nbHasard):\\
\hspace*{2em}\textcolor{magenta}{if} rep<nbHasard:\\
\hspace*{4em}print("+")\\
\hspace*{2em}\textcolor{magenta}{if} rep>nbHasard:\\
\hspace*{4em}print("-")\end{minipage}}@

essai = 11
#on crée une variable qui va stocker le nombre à trouver .
nbHasard = random.randint(1 , 100)
rep = 0

while rep != nbHasard and essai > 0:
    essai = essai - 1
    rep = @\underline{demander\_nombre()}@
    @\underline{aide(rep, nbHasard)}@

if essai > 0 :
    print("Bravo, vous avez trouvé la bonne réponse !")
else :
    print("Vous avez perdu !")
\end{lstlisting}

\methode
\begin{enumerate}
  \item Une \textbf{définition} commence par le mot-clé \code{def}, suivi du nom, des parenthèses et des deux-points. Elle comprend \textbf{tout le corps} : toutes les lignes indentées sous le \code{def}. On l'entoure en entier.
  \item Un \textbf{appel} est le nom d'une fonction suivi de parenthèses, \textbf{sans} \code{def} devant : c'est le moment où la fonction s'exécute.
  \item Une fonction définie mais jamais appelée ne fait rien. Ici, les deux fonctions sont appelées à chaque tour de la boucle \code{while}.
\end{enumerate}
\begin{erreur}
Souligner la ligne \code{def demander\_nombre():} comme un appel. La ligne \code{def} \textbf{définit} la fonction (elle explique ce qu'elle fera) ; elle ne l'exécute pas.
\end{erreur}


\begin{lstlisting}
from turtle import *

def surprise(cote, repetition):
  # Au depart la tortue est a l'horizontale.
  for i in range(repetition):
    forward((i//4+1)*cote) # forward sert a avancer
    right(90) # right sert a tourner a droite

surprise(50, 16)
\end{lstlisting}

\question[2] \enonce{Dessiner ce qu'affiche le code ci-dessus.}
\begin{reponse}
\textbf{Quatre carrés emboîtés}, de côtés 50, 100, 150 et 200 pixels, qui partagent le même coin de départ (en haut à gauche) et s'agrandissent vers la droite et vers le bas.
\begin{center}
\begin{tikzpicture}[scale=0.02]
  \draw[very thick] (0,0) rectangle (50,-50);
  \draw[very thick] (0,0) rectangle (100,-100);
  \draw[very thick] (0,0) rectangle (150,-150);
  \draw[very thick] (0,0) rectangle (200,-200);
  \fill (0,0) circle (4);
  \node[above left, font=\footnotesize] at (0,0) {départ};
  \draw[-{Stealth}, thick] (5,8) -- (45,8);
  \node[above, font=\footnotesize] at (25,10) {1\textsuperscript{er} trait};
  \node[font=\footnotesize] at (25,-25) {50};
  \node[font=\footnotesize] at (75,-75) {100};
  \node[font=\footnotesize] at (125,-125) {150};
  \node[font=\footnotesize] at (175,-175) {200};
\end{tikzpicture}
\end{center}
\end{reponse}
\methode\ Faire le tableau des 16 tours, avec la longueur \code{(i//4+1)*cote} :
\begin{center}
\begin{tabular}{|c|c|c|l|}
\hline
Valeurs de \code{i} & \code{i//4} & Longueur (\code{cote = 50}) & Ce qui est dessiné \\ \hline
0, 1, 2, 3 & 0 & $1 \times 50 = 50$ & un carré de côté 50 \\
4, 5, 6, 7 & 1 & $2 \times 50 = 100$ & un carré de côté 100 \\
8, 9, 10, 11 & 2 & $3 \times 50 = 150$ & un carré de côté 150 \\
12, 13, 14, 15 & 3 & $4 \times 50 = 200$ & un carré de côté 200 \\ \hline
\end{tabular}
\end{center}
\begin{enumerate}
  \item \code{i//4} est le quotient entier de \code{i} par 4 : il augmente de 1 tous les 4 tours.
  \item Pendant 4 tours, la tortue avance de la même longueur et tourne de 90° à droite : elle trace un carré et revient à son point de départ, tournée dans la même direction qu'au début.
  \item Le carré suivant repart donc du même point, avec un côté plus long.
  \item Au départ, la tortue regarde vers la droite ; en tournant à droite, elle descend : les carrés s'étendent vers la droite et vers le bas.
\end{enumerate}
\begin{erreur}
Dessiner une spirale. Il n'y aurait une spirale que si la longueur changeait à \textbf{chaque} trait ; ici elle ne change que tous les 4 traits, après un tour complet.
\end{erreur}

\begin{lstlisting}
def f(x, y, z):
    return x + y, x + z, y + z

def g(x, y):
    if x > 10 :
        return x, y
    else :
        return y, x

def combo(x, y, z):
    s1, s2, s3  = f(x, y, z)
    print(s1, s2, s3)
    a, b = g(s1, s3)
    print(a, b)
    return s1 - a, s2 - b

res1, res2 = combo(4 ,9 ,2)
print(res1, res2)
\end{lstlisting}

\question[3] \enonce{Qu'affiche le code ci-dessus ?}
\begin{reponse}
\begin{lstlisting}[style=sortie]
13 6 11
13 11
0 -5
\end{lstlisting}
\end{reponse}
\methode\ On suit les appels dans l'ordre d'exécution. Seule la ligne 17 lance le calcul : les lignes \code{def} ne font que définir.
\begin{enumerate}
  \item Ligne 17 : appel \code{combo(4, 9, 2)}. Dans \code{combo}, \code{x = 4}, \code{y = 9}, \code{z = 2}.
  \item Ligne 11 : appel \code{f(4, 9, 2)}. \code{f} renvoie trois valeurs : $4 + 9 = 13$, $4 + 2 = 6$, $9 + 2 = 11$. Donc \code{s1 = 13}, \code{s2 = 6}, \code{s3 = 11}. Ligne 12 : affichage \code{13 6 11}.
  \item Ligne 13 : appel \code{g(13, 11)}. Dans \code{g}, \code{x = 13} et \code{y = 11}. $13 > 10$ est vrai : \code{g} renvoie \code{(13, 11)} sans les échanger. Donc \code{a = 13}, \code{b = 11}. Ligne 14 : affichage \code{13 11}.
  \item Ligne 15 : \code{combo} renvoie $13 - 13 = 0$ et $6 - 11 = -5$. Donc \code{res1 = 0}, \code{res2 = -5}. Ligne 18 : affichage \code{0 -5}.
\end{enumerate}
\begin{center}
\begin{tabular}{|l|c|c|c|c|c|c|c|}
\hline
Étape & \code{s1} & \code{s2} & \code{s3} & \code{a} & \code{b} & \code{res1} & \code{res2} \\ \hline
après \code{f} & 13 & 6 & 11 & & & & \\
après \code{g} & 13 & 6 & 11 & 13 & 11 & & \\
après \code{combo} & & & & & & 0 & $-5$ \\ \hline
\end{tabular}
\end{center}
\begin{erreur}
Croire que le \code{x} de \code{g} est le \code{x} de \code{combo} (4). Chaque fonction a ses propres variables locales : dans \code{g}, \code{x} reçoit le premier argument de l'appel, c'est-à-dire \code{s1 = 13}.
\end{erreur}

% =====================================================================
\section{Rédaction de fonction}

\question[2] \enonce{Écrire une fonction \code{division} qui renvoie le quotient de deux nombres donnés en paramètres, en prenant garde que le diviseur soit différent de 0.}
\begin{reponse}
\begin{lstlisting}
def division(a, b):
    if b == 0:
        return None      # division impossible
    return a / b

print(division(7, 2))    # affiche 3.5
print(division(7, 0))    # affiche None
\end{lstlisting}
\end{reponse}
\methode
\begin{enumerate}
  \item L'en-tête : \code{def}, le nom \code{division}, deux paramètres (le dividende \code{a} et le diviseur \code{b}), et les deux-points.
  \item Avant de diviser, tester le cas interdit : si \code{b == 0}, on renvoie \code{None} (la valeur Python qui veut dire « rien »). Un \code{return} arrête la fonction tout de suite.
  \item Sinon, on renvoie \code{a / b} avec \code{return} (et pas \code{print} : la fonction doit \textbf{renvoyer} le résultat).
\end{enumerate}
\textit{Variantes acceptées :} afficher un message d'erreur et renvoyer \code{None} ; utiliser une assertion \code{assert b != 0, "diviseur nul"} ; renvoyer le quotient entier \code{a // b} si « quotient » est compris au sens de la division entière.
\begin{erreur}
Oublier le \code{return} et écrire \code{print(a / b)} : la fonction affiche le résultat mais renvoie \code{None}. \code{x = division(7, 2)} donnerait alors \code{x = None}.
\end{erreur}

\end{questions}

\end{document}
